\section*{अध्याय \ref{ch:b2:fragment-orbitals} --- अणुखंड कक्षक और बहुपरमाणुक अणु}

\begin{solution}{exo:b2:fragment-orbitals:1}
$h_1 + h_2$ और $h_1 - h_2$। तल $xz$, जो कोण को समद्विभाजित करता है और
अणु के लंबवत है, दोनों हाइड्रोजनों का आदान-प्रदान करता है: $h_1 - h_2$ चिह्न
बदलता है, $h_1 + h_2$ नहीं।
\end{solution}

\begin{solution}{exo:b2:fragment-orbitals:2}
छह संयोजकता कक्षक (ऑक्सीजन का 2s और तीन 2p, दो हाइड्रोजन 1s), आठ संयोजकता
इलेक्ट्रॉन, चार भरे कक्षक।
\end{solution}

\begin{solution}{exo:b2:fragment-orbitals:3}
भरे संयोजकता स्तर एक $a_1$ कक्षक और तीन अपभ्रष्ट $t_2$
कक्षकों का समुच्चय हैं: दो ऊर्जाएँ, दो बैंड। $t_2$ बैंड छह इलेक्ट्रॉन रखने वाले तीन
कक्षकों से इलेक्ट्रॉन हटाता है, जबकि $a_1$ के लिए एक कक्षक और दो इलेक्ट्रॉन।
\end{solution}

\begin{solution}{exo:b2:fragment-orbitals:4}
$\pi$ आबंध के लिए कार्बनों के दोनों p कक्षकों का समांतर होना आवश्यक है, जो
दोनों \ce{CH2} तलों को एक उभयनिष्ठ तल में रख देता है।
\end{solution}

\begin{solution}{exo:b2:fragment-orbitals:5}
अमोनिया अपने एकाकी युग्म से इलेक्ट्रॉन खोती है, जो नाइट्रोजन पर अधिकांशतः \omterm{def:b2:diatomic-mos:bonding}{अनाबंधी कक्षक} है;
मेथेन एक आबंधी $t_2$ कक्षक से। अनाबंधी इलेक्ट्रॉन आबंधी इलेक्ट्रॉन से ऊपर,
परमाण्वीय स्तर के अधिक निकट होता है: वह अधिक आसानी से हटाया जाता है।
\end{solution}

\begin{solution}{exo:b2:fragment-orbitals:6}
रैखिक जल के दो अपभ्रष्ट भरे अनाबंधी p कक्षक होते हैं; मोड़ना उनमें से एक को नीचे
लाता है और एक आबंधी स्तर को थोड़ा ऊपर उठाता है: आठ इलेक्ट्रॉनों के साथ शुद्ध लाभ। \ce{BeH2}
के केवल चार हैं, जो दोनों \omterm{def:b2:diatomic-mos:bonding}{आबंधी कक्षकों} को भरते हैं; जिस कक्षक को मोड़ना
नीचे लाता, वह खाली है, जबकि आबंधी स्तर अतिव्यापन खोते हैं: रैखिक आकृति सर्वोत्तम है।
\end{solution}

\begin{solution}{exo:b2:fragment-orbitals:7}
धनायनी केंद्र के पास वाले कार्बन पर हर C--H या C--C आबंध
खाली p कक्षक के साथ संरेखित हो सकता है और लगभग $2\beta^2/\Delta$ स्थायीकरण देता है। \ce{CH3+} का ऐसा कोई
पड़ोसी नहीं; एथिल, आइसोप्रोपिल और tert-ब्यूटिल धनायनों के केंद्र के पास एक, दो और तीन मेथिल
समूह हैं, हर एक एक संरेखित आबंध प्रदान करता है: स्थायित्व उसी
क्रम में बढ़ता है।
\end{solution}

\begin{solution}{exo:b2:fragment-orbitals:8}
दो समान स्तरों के लिए विपाटन $2|\beta|$ है और दोनों $\pi$ इलेक्ट्रॉन $2|\beta|$ पाते हैं,
जो $\cos\theta$ के समानुपाती है। यह $\theta = 60^\circ$ पर आधा और
$90^\circ$ पर शून्य हो जाता है।
\end{solution}

\begin{solution}{exo:b2:fragment-orbitals:9}
बोरेन के छह संयोजकता इलेक्ट्रॉन हैं: वे तीनों \omterm{def:b2:diatomic-mos:bonding}{आबंधी कक्षकों} को भरते हैं, और तल के लंबवत
बोरॉन का p कक्षक खाली रहता है; पिरामिडी बनने से कुछ लाभ नहीं होता।
अमोनिया के दो और हैं: वे उस कक्षक में रहते हैं जो समतलीय आकृति में
अनाबंधी है, और, जल की तरह, पिरामिडी बनना उसे 2s कक्षक और
हाइड्रोजन संयोजन के साथ मिश्रित होकर नीचे गिरने देता है।
\end{solution}

\begin{solution}{exo:b2:fragment-orbitals:10}
$c_1 = c_2 = c_3$ के लिए सेकुलर समीकरण $\alpha + 2\beta$ देते हैं; उसके लांबिक दोनों संयोजन
$\alpha - \beta$ (अपभ्रष्ट) देते हैं। $\alpha + 2\beta$ में दो इलेक्ट्रॉन: ऊर्जा
$2\alpha + 4\beta$, जो $\ce{H2} + \ce{H+}$ ($2\alpha + 2\beta$) से नीचे है क्योंकि $\beta < 0$।
\end{solution}

\begin{solution}{exo:b2:fragment-orbitals:11}
$\pi_1$: तीनों गुणांक एक ही चिह्न के, कोई नोड नहीं (आबंधी); $\pi_2$: सिरे के कार्बनों पर
विपरीत चिह्न और केंद्र पर एक नोड (अनाबंधी); $\pi_3$: एकांतर
चिह्न, दो नोड (प्रतिआबंधी)। ऋणायन के चार $\pi$ इलेक्ट्रॉन हैं: $\pi_2$ उसका
उच्चतम भरा स्तर है, जिसका घनत्व केवल सिरे के कार्बनों पर है।
\end{solution}

\begin{solution}{exo:b2:fragment-orbitals:12}
अक्ष को समाविष्ट करने वाले दोनों तलों में से हर एक में तीन समांतर p कक्षक (O, C, O)
एक आबंधी, एक अनाबंधी (केवल ऑक्सीजनों पर) और एक \omterm{def:b2:diatomic-mos:bonding}{प्रतिआबंधी कक्षक} देते हैं: कुल छह
$\pi$ कक्षक। दोनों आबंधी और दोनों अनाबंधी भरे हैं: आठ
$\pi$ इलेक्ट्रॉन (लुइस चित्र में दो $\pi$ आबंध और ऑक्सीजन के दो एकाकी युग्म)।
\end{solution}

\begin{solution}{pb:b2:fragment-orbitals:1}
\textbf{1.} आठ: छह ऑक्सीजन से, एक-एक हर हाइड्रोजन से।
\textbf{2.} $h_1 + h_2$ और $h_1 - h_2$।
\textbf{3.} 2s और $2\mathrm p_z$।
\textbf{4.} $2\mathrm p_y$।
\textbf{5.} $2\mathrm p_x$, आण्विक तल के लंबवत।
\textbf{6.} छह; चार भरे।
\textbf{7.} $h_1 - h_2$, जो $2\mathrm p_z$ की तरह, ऑक्सीजन पर अक्ष के लंबवत तल से होकर
चिह्न बदलता है।
\textbf{8.} $h_1 + h_2$।
\textbf{9.} $2\mathrm p_x$ और $2\mathrm p_y$, अक्ष के लंबवत: एक अपभ्रष्ट
अनाबंधी युग्म।
\textbf{10.} (2s, $h_1 + h_2$) आबंधी$^2$, ($2\mathrm p_z$, $h_1 - h_2$) आबंधी$^2$, फिर
$2\mathrm p_x^2\,2\mathrm p_y^2$।
\textbf{11.} नहीं: अपभ्रष्ट युग्म भरा है, हर इलेक्ट्रॉन युग्मित है।
\textbf{12.} ऑक्सीजन 2p कक्षक की ऊर्जा पर: वह अनाबंधी है।
\textbf{13.} $2\mathrm p_y$ (रैखिक अणु का अक्ष $z$ के अनुदिश, मोड़ना $yz$ में)।
\textbf{14.} वह 2s और $h_1 + h_2$ के साथ मिश्रित होकर नीचे गिरता है: $3a_1$ स्तर।
\textbf{15.} ($2\mathrm p_z$, $h_1 - h_2$) आबंधी स्तर, जो कुछ अतिव्यापन खोता है जब
हाइड्रोजन अक्ष छोड़ते हैं: यह $1b_2$ बन जाता है।
\textbf{16.} नीचे गिरता भरा स्तर उससे अधिक पाता है जितना दूसरा खोता है: आठ
इलेक्ट्रॉनों के साथ मुड़ी आकृति अधिक स्थायी है।
\textbf{17.} $104.48^\circ$, $90^\circ$ (केवल p कक्षकों से बने आबंध) और $109.5^\circ$ के बीच:
2s कक्षक आबंधन में कुछ भाग लेता है।
\textbf{18.} $2\mathrm p_x$, जो अनाबंधी $1b_1$ ही रहता है।
\textbf{19.} चार।
\textbf{20.} $1b_1$ से, जो ऑक्सीजन का अनाबंधी $2\mathrm p_x$ कक्षक है।
\textbf{21.} अनाबंधी इलेक्ट्रॉन हटाना आबंधों को मुश्किल से बदलता है: आयन की
ज्यामिति अणु जैसी रहती है और थोड़ी कंपन ऊर्जा उत्तेजित होती है। आबंधी इलेक्ट्रॉन
हटाना आबंध लंबाइयाँ और कोण बदलता है, और बैंड
अनेक कंपन स्तरों पर फैल जाता है।
\textbf{22.} \qty{12.62}{eV}, जबकि \qty{13.62}{eV}: जल में ऑक्सीजन हाइड्रोजनों से
खींचा गया आंशिक ऋणात्मक आवेश रखती है, जो उसके 2p स्तरों को ऊपर उठाता है।
\textbf{23.} वे दो समान स्तर नहीं हैं: उनके दोनों संयोजन $3a_1$ और
$1b_1$ कक्षक हैं, भिन्न ऊर्जाओं के। दोनों चित्र उसी कुल
इलेक्ट्रॉन घनत्व का वर्णन करते हैं; केवल विस्थानीकृत कक्षक आयनन की ऊर्जाएँ देते हैं।
\textbf{24.} जल के \textbf{चार} संयोजकता आयनन बैंड हैं, सबसे निचला
$\boldsymbol{\approx \qty{12.6}{eV}}$ पर।
\end{solution}
